\documentclass[12pt,a4paper]{article}
\usepackage[T2A]{fontenc}     % Жестко задает внутреннюю кодировку шрифтов для кириллицы
\usepackage[utf8]{inputenc}   % Говорит, что сам файл сохранен в UTF-8
\usepackage[russian,english]{babel} % Включает поддержку переносов и словарей
\usepackage{amsmath}
\usepackage{hyperref}

\title{MZLS-Engine: Architecture Decision Records (ADR)}
\author{kshakirov}
\date{July 2026}

\begin{document}

\maketitle

\section*{ADR-0001: Выбор сетевого ввода-вывода / Network I/O Core}
\textbf{Status:} Accepted \\

\subsection*{RU: Контекст и Обоснование}
Для создания веб-сервера МЗ с нуля выбран базис блокирующего ввода-вывода (BIO) на базе Unix Domain Sockets (UDS), изолированный за обратным прокси Nginx. 
Эмпирические тесты по методу Фрэнсиса Бэкона на 8-летнем аппаратном обеспечении HP выявили предел производительности системы. Плотность нагрузки $12$ потоков / $100$ коннектов приводит к Context Switching ядра ОС, тогда как балансировка на $6$ потоков / $40$ коннектов обеспечивает стабильные $23\,466$ RPS с математическим нулем ошибок ($0\%$ Non-2xx).

\subsection*{EN: Context \& Decision}
The core web server architecture uses Blocking I/O (BIO) over Unix Domain Sockets (UDS) behind an Nginx reverse proxy. 
Empirical stress testing (Francis Bacon method) on 8-year-old HP hardware defined the system's operational threshold. High-density traffic ($12$ threads / $100$ connections) causes severe OS context switching, whereas balanced traffic ($6$ threads / $40$ connections) delivers a stable $23,466$ RPS with a mathematical zero error rate ($0\%$ Non-2xx).




\section*{ADR-0002: Инкрементальный мониторинг латентности по алгоритму Велфорда / Incremental Latency Monitoring via Welford's Algorithm}
\textbf{Status:} Proposed \\

\subsection*{RU: Контекст и Математическое Обоснование}
Для сбора метрик задержки (latency) веб-сервера МЗ под нагрузкой недопустимо использовать массивы данных или блокирующие примитивы синхронизации, так как это вызывает аллокационный кризис в Heap и деградацию кэша процессора. 

Расчет среднего ($\mu$) и дисперсии ($\sigma^2$) выполняется инкрементально за строгое время $O(1)$ по формулам Велфорда:
$$\mu_n = \mu_{n-1} + \frac{x_n - \mu_{n-1}}{n}$$
$$M_{2,n} = M_{2,n-1} + (x_n - \mu_{n-1})(x_n - \mu_n)$$
$$\sigma^2 = \frac{M_{2,n}}{n}$$

Для защиты от ложного разделения кэш-линий (False Sharing) в архитектуре x86 процессора HP, класс \texttt{WelfordState} проектируется с жестким выравниванием (padding) по границе $64$ байт с помощью пустых полей-заполнителей типа \texttt{long}. Потоки пула пишут метрики асинхронно в свои изолированные локальные состояния без блокировок, а фоновый поток-наблюдатель производит консолидацию данных (Reduce).

\subsection*{EN: Context \& Mathematical Specification}
To collect latency metrics under high load, capturing samples into arrays or using blocking synchronization primitives is prohibited due to memory allocation overhead in the Heap and CPU cache degradation.

The running mean ($\mu$) and variance ($\sigma^2$) are calculated incrementally in $O(1)$ time using Welford's equations as stated above. To prevent False Sharing across the x86 CPU cache lines on the host HP hardware, the \texttt{WelfordState} structure enforces a strict $64$-byte cache-line alignment using dummy \texttt{long} padding fields. Worker threads update their isolated thread-local states lock-free, while a background observer thread executes a periodic asynchronous Reduce operation.


\section*{ADR-0003: Детерминированный конечный автомат парсера HTTP / Deterministic Finite Automaton for Zero-Allocation HTTP Parsing}
\textbf{Status:} Proposed \\

\subsection*{RU: Математическая модель Автомата}
Парсер HTTP-заголовков проектируется как строго детерминированный конечный автомат (ДКА), работающий по алфавиту ASCII-байтов за время $O(n)$ и с памятью $O(1)$. Память в Heap не аллоцируется. На входе автомат принимает сырой буфер чтения. 

Множество состояний автомата $S = \{ \text{START}, \text{METHOD}, \text{URI}, \text{VERSION}, \text{H\_START}, \text{H\_NAME}, \text{H\_VALUE}, \text{FINISHED} \}$.
Таблица переходов по критическим триггерам (байтам):
\begin{itemize}
    \item $\text{START} \xrightarrow{\text{ASCII char}} \text{METHOD}$
    \item $\text{METHOD} \xrightarrow{\text{0x20 (Space)}} \text{URI}$ [Фиксация координат метода]
    \item $\text{URI} \xrightarrow{\text{0x20 (Space)}} \text{VERSION}$ [Фиксация координат URI]
    \item $\text{VERSION} \xrightarrow{\text{0x0A (\textbackslash n)}} \text{H\_START}$
    \item $\text{H\_START} \xrightarrow{\text{0x0D (\textbackslash r)}} \text{проверка на 0x0A} \rightarrow \text{FINISHED}$ [Конец заголовков HTTP]
    \item $\text{H\_START} \xrightarrow{\text{ASCII char}} \text{H\_NAME}$
    \item $\text{H\_NAME} \xrightarrow{\text{0x3A (:)}} \text{H\_VALUE}$
    \item $\text{H\_VALUE} \xrightarrow{\text{0x0A (\textbackslash n)}} \text{H\_START}$
\end{itemize}

Выходом автомата является структура метаданных (только примитивы \texttt{int}), хранящая исключительно \texttt{offset} и \texttt{length} для каждого токена относительно исходного буфера сокета.

\subsection*{EN: DFA Formal Specification}
The HTTP header parser is designed as a strict Deterministic Finite Automaton (DFA) operating over the ASCII byte alphabet with $O(n)$ time complexity and $O(1)$ space complexity. Zero memory is allocated in the Java Heap.

The automaton keeps track of token boundaries using integer pairs (\texttt{offset} and \texttt{length}) referencing the underlying raw socket buffer. The parser halts execution immediately upon reaching the \text{FINISHED} state, preventing any lookahead or reading into the request HTTP body payload.



\end{document}
